\subsection{Synchrotron Wiggler Emission: a summary}

Wigglers and short insertion devices (IDs) are increasingly used as photon sources in modern synchrotron facilities, particularly in upgraded storage rings like ESRF-EBS. Unlike undulators, wiggler radiation is \textbf{fully incoherent}, making existing undulator calculation methods inefficient. The model used in SHADOW4 is summarized here.

\subsubsection{Key Assumptions}

The method assumes:
\begin{enumerate}
    \item \textbf{Fully incoherent emission}: Radiation from different points along the electron trajectory superimposes incoherently.
    \item \textbf{Bending magnet approximation}: Emission at each trajectory point equals that of a bending magnet with the same local curvature.
\end{enumerate}

\subsubsection{Calculation Method}

\textbf{Electron Trajectory Calculation}

%Given a tabulated vertical magnetic field $B_y(z)$, the electron trajectory is computed:

The wiggler magnetic field for a
conventional wiggler has only a vertical component $(0, 0, B_z)$. Suppose we know a tabulation of $B_z(y)$ (measured or calculated, e.f. sinusoilal).  An
electron entering in this magnetic field will have a transversal velocity given by the integral of the magnetic field


\begin{equation}
\beta_x(y) = -\frac{c \times 10^{-9}}{E_e[GeV]} \int_{y_0}^{y} B_z(s) \, ds.
\end{equation}

The electron trajectory is in the horizontal plane and is given
by the integral of the velocity

\begin{equation}
x(y) = \int_{y_0}^{y} \beta_x(s) \, ds.
\end{equation}

Note that in both integrals an additional constant can be added in order to obtain ``centered" profiles. This is discussed in the practical exercise, and it will be whown that an adequate selection of the constant is important to get the radiation centered on axis. 

\textbf{Local Curvature and Critical Energy}

The local curvature (or inverse of the local radius) is also
calculated, because the local emission corresponds to that of
a bending magnet of that curvature. It is proportional to the
second derivative of the trajectory, therefore proportional to $B_z$

\begin{equation}
R^{-1}(y) = \frac{e}{\gamma m_e c} B_z(y)
\end{equation}

with $e$ the electron charge, $m_e$ the electron mass and $\gamma$ the
electron energy in units of $m_e c^2$. Each curvature has associated a bending magnet critical energy given by

\begin{equation}
E_c(y) = \frac{3}{4\pi} \lambda \gamma^3 R^{-1}(y)
\end{equation}

\textbf{Wiggler Spectrum Calculation}

The total spectrum $F(E)$ is obtained by:
\begin{enumerate}
    \item Calculating bending magnet emission at each $z$ point
    \item Multiplying by $(R(z)10^{-3})^{-1}$ for arc length conversion
    \item Accounting for trajectory projection via $\sqrt{1 + (\beta_x(z)/\beta_z(z))^2}$
    \item Integrating over $z$
\end{enumerate}


\subsubsection{Sampling mechanism}

The rays are obtained by:
\begin{enumerate}
    \item Sample the energy
    \item Sample the position in the trajectory
    \item Add electron emittance
\end{enumerate}

\subsubsection{Conclusion}

The presented algorithm provides a fast, accurate method for ray tracing simulations of wigglers and short insertion devices. It efficiently handles the incoherent nature of wiggler radiation and accommodates non-sinusoidal magnetic fields.
